\documentclass[10pt,a4paper]{amsart}
\usepackage[margin=19mm]{geometry}
\usepackage{amsmath,amssymb,mathtools}
\usepackage[scr=rsfs,cal=euler]{mathalfa}
\usepackage{xfrac}
\usepackage[unicode,hidelinks,bookmarks=false]{hyperref}
\newcommand{\ie}{i.e.\ }
\newcommand{\cf}{cf.\ }
\newcommand{\resp}{respectively\ }
\newcommand{\Subsection}{\subsection}
\providecommand{\nobreakdash}{\nobreak\hbox{-}}
\setlength{\emergencystretch}{2em}
\multlinegap0pt
\usepackage{bm}
\usepackage{bbm}
\usepackage{dsfont}
\usepackage{xspace}
\usepackage{cancel}
\usepackage{mathtools}
\usepackage{amsthm}
\makeatletter
\@ifundefined{theorem}{\newtheorem{theorem}{Theorem}}{}
\@ifundefined{definition}{\newtheorem{definition}[theorem]{Definition}}{}
\@ifundefined{lemma}{\newtheorem{lemma}[theorem]{Lemma}}{}
\makeatother
\def\bN{\mathbb{N}}
\def\bR{\mathbb{R}}\def\bS{\mathbb{S}}
\def\falpha{{\boldsymbol\alpha}}
\def\fbeta{{\boldsymbol\beta}}
\def\fgamma{{\boldsymbol\gamma}}
\def\fk{\mathbf k}\def\fl{\mathbf l}
\def\flambda{{\boldsymbol\lambda}}
\def\fl{\mathbf l}
\def\fn{\mathbf n}\def\fN{\mathbf N}
\def\fone{{\boldsymbol 1}}
\def\foo{{\boldsymbol 0}}
\def\ft{\mathbf t}
\def\fx{\mathbf x}
\def\fy{\mathbf y}\def\fz{\mathbf z}
\def\fz{\mathbf z}
\newcommand{\Id}{\mbox{Id}}
\newcommand{\abs}[1]{\left\vert #1 \right\vert}
\newcommand{\compAbs}[1]{\left\vert#1\right\vert_{\text{c}}}
\newcommand{\norm}[1]{\left\Vert #1\right\Vert}
\title[A Stochastic Reconstruction Theorem on Rectangular Increment]{A Stochastic Reconstruction Theorem on Rectangular Increments with an Application to a Mixed Hyperbolic SPDE}
\author{Carlo Bellingeri, Hannes Kern}
\date{Source: arXiv:2404.18634v3 (CC BY 4.0)}
\begin{document}
\maketitle

\noindent\emph{Source and licence.} A Stochastic Reconstruction Theorem on Rectangular Increments with an Application to a Mixed Hyperbolic SPDE. Version arXiv:2404.18634v3 (CC BY 4.0), distributed under the Creative Commons Attribution 4.0 International licence, as recorded in the arXiv licence metadata for this exact version and shown on the abstract page https://arxiv.org/abs/2404.18634v3. Full rights notice: licenses/arxiv-2404.18634-CC-BY-4.0.txt.\par\smallskip

\noindent\textit{Editorial scope.} Counted unit: the Lemma whose source label is \texttt{lem:TaylorRecIncrement} in the article (source0.tex lines 2055--2072, source label \texttt{lem:TaylorRecIncrement}), delivered together with the article's own complete original proof of it (source0.tex lines 2074--2139, one \texttt{proof} environment of 66 lines). No mathematics is written, repaired or re-derived: statement and proof are the article's own text line for line. Required context retained uncounted: rectangular\_increment (lines 431--434); eq\_det\_Calpha (lines 819--821); def:Calpha (lines 2010--2016). Every definition, display and label of those lines is retained unchanged. Omitted from this package and not redistributed: 1--430, 435--818, 822--2009, 2017--2054, 2073--2073, 2140--2223. Nothing inside a retained excerpt is omitted or altered: every retained body is delivered exactly as the article's lines are written, so no omission receipt of any kind accompanies this package. Citations: the retained excerpts carry no citation command at all, so no bibliography entry of the article is transcribed and no reference is redistributed. Reference adaptation: none was needed and none was made; every reference inside the delivered unit resolves inside it, so no unresolved reference or citation is produced. Printed slips kept literal: no printed word, symbol, hypothesis or number of the counted statement or of its proof was altered. Layout and class: the article's own class is \texttt{article} with packages including inputenc, babel, amsmath, amssymb, mathrsfs, color, amsthm, mathtools, bm, enumerate, hyperref, fullpage, graphicx, caption; the wrapper is the lane's pinned \texttt{amsart} editorial wrapper at 10pt on A4, which restates the theorem-like environments the retained text needs and adds the article's own macro definitions that the retained text needs (\texttt{\textbackslash Id}, \texttt{\textbackslash abs}, \texttt{\textbackslash bN}, \texttt{\textbackslash bR}, \texttt{\textbackslash compAbs}, \texttt{\textbackslash falpha}, \texttt{\textbackslash fbeta}, \texttt{\textbackslash fgamma}, \texttt{\textbackslash fk}, \texttt{\textbackslash fl}, \texttt{\textbackslash flambda}, \texttt{\textbackslash fn}, \texttt{\textbackslash fone}, \texttt{\textbackslash foo}, \texttt{\textbackslash ft}, \texttt{\textbackslash fx}, \texttt{\textbackslash fy}, \texttt{\textbackslash fz}, \texttt{\textbackslash norm}), with extra wrapper packages (bm, bbm, dsfont, xspace, cancel, mathtools). The delivered numbering is the unit's own. Assets: no retained span contains a figure environment, a graphics-inclusion command, a PDF-page inclusion, a table or a TikZ picture; no image file is copied into this package, the package declares no asset dependency and no image file is redistributed with it. Licence: version arXiv:2404.18634v3 of this article is distributed under the Creative Commons Attribution 4.0 International licence, as recorded in the arXiv licence metadata for this exact version and shown on the abstract page https://arxiv.org/abs/2404.18634v3. A full rights notice is delivered at licenses/arxiv-2404.18634-CC-BY-4.0.txt. Characters: every retained excerpt is pure ASCII, so the delivered engine, XeLaTeX with its default Latin Modern text font, reports no missing character.
\begin{align}\label{rectangular_increment}
\square_{\fx,\fy}^\theta f :=&\prod_{i\in\theta}(\text{Id} - \pi^i_{\fx})f(\pi^{\theta^c}_{\fx}\fy) = \sum_{\theta'\subset\theta} (-1)^{\sharp(\theta\setminus\theta')} f(\pi^{\theta'}_{\fy}\fx)
= \prod_{i\in\theta}(\pi^i_\fy - \Id)f(\fx)\,, 
\end{align}

\begin{equation}\label{eq_det_Calpha}
\norm{Y}_{C^\falpha} := \sum_{\theta\subset[d]} \norm{\square^\theta Y}_{\falpha,\theta} <\infty\,,\quad  \norm{\square^\theta Y}_{\falpha,\theta}:=\sup_{\fx, \fy\in [0,T]^d,\, \fx\neq \fy} \frac{|\square^\theta_{\fx,\fy} Y|}{|\fx- \fy|^{\falpha}_{\text{c}, \theta}}\,.
\end{equation}

\begin{definition}\label{def:Calpha}
    Let $\fbeta\in[0,\infty)^d$ and choose $\fn\in\bN^d, \fgamma\in(0,1]^d$ such that $\fbeta = \fn+\fgamma$. If $\beta_i = 0$ for some $i\in[d]$, we choose $n_i = \gamma_i= 0$. Let $g:[0,T]^d\to \bR$ be such that all derivatives $\partial^\fk g$ exist for $\foo\le\fk\le\fn$. Then we say that $g\in C^\fbeta$ if the norm
    \[
        \norm{g}_{C^\fbeta} := \sum_{\fk\le\fn} \left(\norm{\partial^\fk g}_\infty + \norm{\partial^{\fk} g}_{C^\fgamma}\right)
    \]
    is finite. Here, $\norm{\cdot}_{C^\fgamma}$ is defined as in \eqref{eq_det_Calpha}.
\end{definition}

\begin{lemma}\label{lem:TaylorRecIncrement}
    Let $g\in C^{\fbeta}, \fbeta>\foo$ and let $\fn,\fgamma$ be as in Definition \ref{def:Calpha}. Let $\theta\subset[d], \fx,\fy,\fz\in[0,T]^d$. Then for all $\fk\le\fn$, there exists $A(\fx,\fy,\fk,\theta)\in\bR$ such that
    \begin{align}
    \begin{split}\label{ineq:Taylor1}
        \square^\theta_{\fx,\fy} T^{\fn} g(\fz) &= \sum_{\fk\le\fn} (\fz-\fy)^\fk A(\fx,\fy,\fk,\theta)\\
        \abs{A(\fx,\fy,\fk,\theta)} &\lesssim \compAbs{\fx-\fy}^{\fbeta_\theta-\fk}\norm{g}_{C^\fbeta}\,.
        \end{split}
    \end{align}
    Here, the rectangular increment is taken with respect to the map $\fx\mapsto T^\fn_\fx g(\fz)$.
    
    Furthermore, if $g\in C^{r+\abs{\fn_\theta+\fone_\theta}}$, then for each $\theta\subset[d],~\fx,\fy,\fz\in[0,T]^d$ such that $\fx_\theta = \fz_\theta$ and $\flambda\in(0,1]^d$, there exists a $B(\fy,\fz,\flambda,\theta)\in\bR$, such that
    \begin{align}
    \begin{split}\label{ineq:Taylor2}
        R^\theta g(\fx,\flambda*\fy+\fz,\fn) &= \flambda^{\fn_\theta+1_\theta} B(\fy,\fz,\flambda,\theta)\,, \\
        \norm{B(\cdot,\fz,\flambda,\theta)}_{C^r}&\lesssim \norm{g}_{C^{r+\abs{\fn_\theta+\fone_\theta}}}
    \end{split}
    \end{align}
\end{lemma}

\begin{proof}Let us start with proofing \eqref{ineq:Taylor1}. Writing out the rectangular increment using \eqref{rectangular_increment}, one gets
    \begin{align*}
        \square_{\fx,\fy}^\theta T^\fn g(\fz) &= \left[\prod_{i\in\theta} (\pi^i_{y_i} -\Id) \prod_{i=1}^d T_{x_i}^{i,n_i}\right]g(\fz) \\
        &=\left[\prod_{i\in\theta} (T_{y_i}^{i,n_i}-T_{x_i}^{i,n_i}) \prod_{i\notin\theta} T_{x_i}^{i,n_i}\right]g(\fz)\,.
    \end{align*}
    We write $\tilde g(\fz) = T^{\theta^c,\fn}_\fx g(\fz)$. Observe that writing out the Taylor expansion $T^{i,n_i}_{y_i}$ for some $i\in\theta$ and expanding all $\partial^{i,k_i}g(\pi^i_{y_i}\fz)$ in $x_i$ leads to
     \begin{align*}
        T^{i,n_i}_{y_i} \tilde g(\fz) &= \sum_{k_i=0}^{n_i} \frac{1}{k_i!} \partial^{i,k_i} \tilde g(\pi_{y_i}^i\fz) (z_i-y_i)^{k_i}\\
        &= \sum_{k_i=0}^{n_i}  \frac{(z_i-y_i)^{k_i}}{k_i!} \left( \sum_{l_i= 0}^{n_i-k_i}\frac{1}{l_i!} \partial^{i,k_i+l_i} \tilde g(\pi^i_{x_i}\fz)(y_i-x_i)^{l_i} + R^i\tilde g(x_i,\pi_{y_i}^i\fz,n_i-k_i)\right)\,.
    \end{align*}
    Using the binomial formula leads to
    \begin{equation}\label{eq:NewSupportPoint}
    (T_{y_i}^{i,n_i} -T_{x_i}^{i,n_i})\tilde g(\fz) = \sum_{k_i=0}^{n_i} \frac{(z_i-y_i)^{k_i}}{k_i!} R^i\tilde g(x_i,\pi_{y_i}^i\fz,n_i-k_i)\,.
    \end{equation}
    Note that due to the definition of $R^i\tilde g$, we can write for any $j\neq i$
    \[
    (T^{j,n_j}_{y_j}- T^{j,n_j}_{x_j}) R^i\tilde g(x_i,\pi^i_{y_i}\fz,n_i) = \int_{x_i}^{y_i} \frac{(y_i-t_i)^{n_i-1}}{(n_i-1)!} (\pi^i_{t_i} - \pi^i_{x_i})(T^{j,n_j}_{y_j}- T^{j,n_j}_{x_j})\left(\partial^{i,n_i}\tilde g\right)(\fz) dt_i\,,
    \]
    allowing us to iteratively apply the identity \eqref{eq:NewSupportPoint} to all $i\in\theta$. This leads to the formula
    \begin{equation}\label{eq:rectangIncrementTaylorExpansion}
        \square_{\fx,\fy}^\theta T^\fn g(\fz) = \prod_{i\in\theta}(T^{i,n_i}_{y_i}- T^{i,n_i}_{x_i})\tilde g(\fz) = \sum_{\fk_\theta = \foo}^{\fn_\theta}\prod_{i\in\theta}\frac{(z_i-y_i)^{k_i}}{k_i!} R^\theta \tilde g(\fx,\pi^\theta_\fy \fz,\fn-\fk)\,.
    \end{equation}
    The term $R^\theta \tilde g(\fx,\pi^\theta_\fy \fz,\fn-\fk)$ still depends on $z_i$ for $i\notin\theta$, so we extend it as follows:
    \begin{align}\begin{split}\label{eq:ExpandingGTaylor}
        R^\theta T_\fx^{\theta^c,\fn} g(\fx,\pi^\theta_\fy \fz,\fn-\fk) &= \sum_{\fk_{\theta^c}\le\fn_{\theta^c}} \prod_{i\in\theta^c} \frac{(z_i-x_i)^{k_i}}{k_i!} \\&\qquad\int_{\fx_{\tilde\theta}}^{\fy_{\tilde\theta}} \prod_{i\in\tilde\theta} \frac{(y_i-t_i)^{n_i-k_i-1}}{(n_i-k_i-1)!} \square^\theta_{\pi_\fx^\theta\fz,\pi_{\ft_{\tilde\theta}}^{\tilde\theta}\fy} \partial^{\fn_\theta-\fk_{\theta}+\fk_{\theta^c}}g(\pi_\fx^{\theta^c}\cdot)d\ft_{\tilde\theta}\,,
    \end{split}\end{align}
    where $\tilde\theta\subset\theta$ such that $i\in\tilde\theta$ if and only if $n_i-k_i\ge 1$. Note that
    \begin{align*}
        A_1(\fx,\fy,\fk,\theta) :&= \int_{\fx_{\tilde\theta}}^{\fy_{\tilde\theta}} \prod_{i\in\tilde\theta} \frac{(y_i-t_i)^{n_i-k_i-1}}{(n_i-k_i-1)!} \square^\theta_{\pi_\fx^\theta\fz,\pi_{\ft_{\tilde\theta}}^{\tilde\theta}\fy} \partial^{\fn_\theta-\fk_\theta+\fk_{\theta^c}}g(\pi_\fx^{\theta^c}\cdot)d\ft_{\tilde\theta}\\
        &=\int_{\fx_{\tilde\theta}}^{\fy_{\tilde\theta}} \prod_{i\in\tilde\theta} \frac{(y_i-t_i)^{n_i-k_i-1}}{(n_i-k_i-1)!} \square^\theta_{\fx,\pi_{\ft_{\tilde\theta}}^{\tilde\theta}\fy} \partial^{\fn_\theta-\fk_\theta+\fk_{\theta^c}}gd\ft_{\tilde\theta}
    \end{align*}
    does not depend on $\fz$ at all. Using the definition of $\norm{g}_{C^\fbeta}$, it is easy to see
    \[
        \abs{A_1(\fx,\fy,\fk,\theta)} \lesssim \abs{\fx-\fy}^{\fbeta-\fk}_{\text{c},\theta} \norm{g}_{C^\fbeta}\,.
    \]
    Putting this back into \eqref{eq:ExpandingGTaylor} and using the binomial formula on $(z_i-x_i)^{k_i}, i\in\theta^c$ leads to
    \begin{equation}\label{def_A_2}
        R^\theta T_\fx^{\theta^c,\fn} g(\fx,\pi^\theta_\fy \fz,\fn-\fk) = \sum_{\fl_{\theta^c}\le\fk_{\theta^c}\le\fn_{\theta^c}} (\fz-\fy)^{\fl_{\theta^c}}\underbrace{\prod_{i\in\theta^c} \frac{(y_i-x_i)^{k_i-l_i}}{l_i! (k_i-l_i)!}A_1(\fx,\fy,\fk,\theta)}_{A_2(\fx,\fy,\fk,\fl_{\theta^c},\theta)}\,,
    \end{equation}
    where $A_2$ clearly fulfils
    \[
    \abs{A_2(\fx,\fy,\fk,\fl_{\theta^c},\theta)} \lesssim \abs{\fx-\fy}^{\fbeta_\theta-\fk_\theta+\fk_{\theta^c}-\fl_{\theta^c}}_{\text{c}} \norm{g}_{C^\fbeta}\,.
    \]
    By summing out $\fk_{\theta^c}$, we get
    \begin{align}\begin{split}\label{def_A_3}
        R^\theta T_\fx^{\theta^c,\fn} g(\fx,\pi^\theta_\fy \fz,\fn-\fk) &= \sum_{\fl_{\theta^c}\le\fn_{\theta^c}} (\fz-\fy)^{\fl_{\theta^c}}\underbrace{\sum_{\fl_{\theta^c}\le\fk_{\theta^c}\le\fn_{\theta^c}}A_2(\fx,\fy,\fk,\fl_{\theta^c},\theta)}_{A_3(\fx,\fy,\fk_{\theta},\fl_{\theta^c},\theta)}\\
        \abs{A_3(\fx,\fy,\fk_\theta,\fl_{\theta^c},\theta)} &\lesssim \abs{\fx-\fy}^{\fbeta_\theta-\fk_\theta-\fl_{\theta^c}}_{\text{c}} \norm{g}_{C^\fbeta}
    \end{split}\end{align}
    Here, $\sum_{\fl_{\theta^c}\le\fk_{\theta^c}\le\fn_{\theta^c}}$ is seen as a sum over $\fk_{\theta^c}$ with fixed $\fl_{\theta^c}$. Putting \eqref{def_A_3} back into \eqref{eq:rectangIncrementTaylorExpansion} with $A(\fx,\fy,\fk,\theta) := A_3(\fx,\fy,\fk_\theta,\fk_{\theta^c},\theta)$ gives \eqref{ineq:Taylor1}.

    We now pass on to prove \eqref{ineq:Taylor2}. If $g$ has more than $n_i$ derivatives in some direction $i\in[d]$, one can work with the easier remainder formula 
    \begin{equation}
        R^ig(x_i,\fz,n_i) =\int_{x_i}^{z_i} \frac{(z_i-t_i)^{n_i}}{n_i!} \partial^{i,n_i+1} g(\pi^i_{t_i}\fz) dt_i\,,
    \end{equation}
    leading to
    \begin{equation}
        R^\theta g(\fx,\fz,\fn) = \int_{\fx_\theta}^{\fz_\theta} \prod_{i\in\theta}\frac{(z_i-t_i)^{n_i}}{n_i!} \partial^{\fn_\theta+\fone_\theta} g(\pi_{\ft_\theta}^\theta \fz) d\ft_\theta\,.
    \end{equation}
    The rest of the proof follows by using the substitution $\ft = \flambda*\tilde\ft + \fx$ and calculating (recall $\fx_\theta = \fz_\theta$)
    \begin{align*}
    R^\theta g(\fx,\flambda*\fy+\fz,\fn) &= \int_{\fx_\theta}^{\fy_\theta*\flambda_\theta +\fx_\theta}\prod_{i\in\theta}\frac{(\lambda_i y_i -x_i-t_i)^{n_i}}{n_i!} \partial^{\fn_\theta+\fone_\theta} g(\pi^\theta_{\ft_\theta} (\flambda*\fy+\fz)) d\ft_\theta\\
    &= \flambda^{\fn_\theta+\fone_\theta} \int_\foo^{\fy_\theta} \prod_{i\in\theta}\frac{(y_i-\tilde t_i)^{n_i}}{n_i!} \partial^{\fn_\theta+\fone_\theta} g(\flambda*\pi^\theta_{\tilde\ft}\fy+\fz)d\tilde\ft_\theta\\
    &= \flambda^{\fn_\theta+\fone_\theta} B(\fy,\fz,\flambda,\theta)\,.
\end{align*}
The definition of $B$ immediately gives $\norm{B(\cdot,\fz,\flambda,\theta)}_{C^r} \lesssim \norm{g}_{C^{r+\abs{\fn_\theta+\fone_\theta}}}$.
\end{proof}

\end{document}
